% notations of definitions of unit 132 that this note uses and another note of the unit states
% (made by shared/mkspec.py from shared/notation-132.json; the same TeX as in the note that states them)
\newcommand{\sieveVector}[2]{\Xi(#1;#2)}
\newcommand{\IsCutBessel}[7]{\operatorname{CB}(#1,#2,#3,#4,#5,#6,#7)}
\newcommand{\formConst}[1]{c_{\mathrm{f}}(#1)}
\newcommand{\HasFibreBound}[5]{\operatorname{Fib}_{#2,#4}(#1;#3)\le #5}
\newcommand{\Levels}[3]{\operatorname{Lev}(#1,#2,#3)}
\newcommand{\floorMeasure}[2]{\theta_{#2}(#1)}
\newcommand{\stackMeasure}[1]{\mu^{#1}}
\newcommand{\groupMeasure}[1]{\mathsf{M}(#1)}
\newcommand{\upperMeasure}[2]{\mathsf{M}_{\ge #2}(#1)}
\newcommand{\IsNear}[3]{#3\in\operatorname{Near}_{#2}(#1)}
\newcommand{\orthonormalPart}[1]{\mathsf{O}(#1)}
\newcommand{\closePairs}[1]{\operatorname{Cl}(#1)}
\newcommand{\floorValues}[4]{\operatorname{Val}_{#2}(#1;#3,#4)}
\newcommand{\volumeExponent}[1]{\mathsf{s}(#1)}
\newcommand{\floorWeight}[4]{\mathsf{w}_{#2}(#1;#3,#4)}
\newcommand{\floorModulus}[3]{\mathsf{d}_{#2}(#1;#3)}
\newcommand{\FloorMassLe}[4]{\operatorname{Mass}_{#2}(#1;#4)\le #3}
\newcommand{\upperPairing}[5]{\mathsf{Av}_{#2}(#1;#3)(#4,#5)}
\newcommand{\IsPairingKernel}[4]{(#3,#4)\in\operatorname{PK}_{#2}(#1)}
\newcommand{\invWeight}[3]{\mathsf{w}^{-1}_{#2}(#1;#3)}
\newcommand{\pivotModulus}[5]{\mathsf{Pv}_{#2}(#1;#3)(#4,#5)}
\newcommand{\weightedSquare}[5]{\mathsf{Sq}_{#2}(#1;#3,#4)(#5)}
\newcommand{\weightedSquareSum}[5]{\mathsf{Sm}_{#4}(#1;#2,#3)(#5)}
\newcommand{\selfNorm}[4]{\operatorname{SN}_{#2}(#1;#3,#4)}
\newcommand{\Gauges}[1]{\operatorname{Gau}(#1)}
\newcommand{\positionsOfSize}[2]{\operatorname{Pos}(#1;#2)}
\newcommand{\headDensity}[4]{\mathsf{h}_{#1}(#4;#2,#3)}
\newcommand{\translatePairing}[5]{\langle #1(#4),#2(#3#5)\rangle}
\newcommand{\IsBessel}[2]{#1\in\operatorname{Bes}(#2)}
\newcommand{\IsCut}[1]{#1\in\operatorname{Cut}}
\newcommand{\pairingNorm}[2]{\operatorname{PN}(#1,#2)}
\newcommand{\formIntegrandOne}[6]{\mathsf{I}^{1}_{#1}(#2;#3,#4,#5)(#6)}
\newcommand{\formIntegrand}[7]{\mathsf{I}_{#1}(#2,#3;#4,#5,#6)(#7)}
\newcommand{\floorForm}[5]{\mathcal{I}_{#1}(#2,#3;#4,#5)}
\newcommand{\formModulus}[6]{\mathcal{I}^{\mathrm{mod}}_{#1}(#2,#3;#4,#5,#6)}
\newcommand{\formMajorant}[6]{\mathcal{I}^{\mathrm{maj}}_{#1}(#2,#3;#4,#5,#6)}
\newcommand{\selfNormProduct}[7]{\operatorname{SNP}_{#1}(#7;#2,#3;#4,#5,#6)}
\newcommand{\Sizes}[1]{\operatorname{Siz}(#1)}
\newcommand{\weylRatio}[3]{w_{#1}(#3)^{-1}w_{#1}(#2)}
\newcommand{\vecsOfSize}[1]{\mathrm{Vec}_{#1}}
\newcommand{\HasMatSize}[2]{\operatorname{sz}(#2)\le #1}
\newcommand{\IsHead}[2]{#2\in\operatorname{Hd}_{#1}}
\newcommand{\floorConst}[2]{C_{#2}(#1)}
\newcommand{\globalConst}[1]{C_{\mathrm{gl}}(#1)}
